% Solutions to the Chapter 6 exercises that are not code, from lectures/L06/appendix/c_exercises.md.
% Exercises 6.7 and 6.8 are Code and are not answered; see chapter.tex for why.

\section{Chapter 6: Memory, MMIO, and Resource Management}
\label{sol:sec:c6}
Answers to the exercises in \secref{c6:sec:exercises}. Exercises~\ref{c6:ex:7} and~\ref{c6:ex:8}
are Code, and are checked on the target as their Check yourself lines say.

% ------------------------------------------------------------------------------------------
\solution[fillaccent2]{c6:ex:1}{Two kinds of address}{Recall}

\xpart{a} \code{void *}: yes; it is a kernel virtual address, such as \code{kmalloc} returns.
\code{phys_addr_t}: no; it is a number, where something sits on the bus, and the CPU can reach it
only through a mapping, which for a device's registers is \code{ioremap}. \code{void __iomem *}: not
directly; only through the accessors, \code{readl} and \code{writel} and their family.
\code{char __user *}: no; only through \code{copy_to_user}, \code{copy_from_user},
\code{get_user} and \code{put_user}, which check the address and survive a fault on it.

\xpart{b} They put a pointer in an address space of its own. When \code{sparse} runs,
\code{__iomem} expands to \code{__attribute__((noderef, address_space(__iomem)))}, and
\code{__user} to the same with its own space (\file{include/linux/compiler_types.h}); for the
compiler they expand to nothing. \code{sparse}, run with \code{make C=1}, then warns wherever such a
pointer is dereferenced or passed where a pointer to another address space is expected: a user
pointer handed to \code{memcpy}, or an \code{__iomem} pointer stored in a \code{void *}. The
annotation records in the type which side owns the address, and the checker enforces it.

\xpart{c} No to both. An entry means that somebody registered the range in the kernel's resource
tree: a driver that called \code{request_mem_region}, or the kernel itself for System RAM and its
own image. It is bookkeeping, and says nothing about whether a virtual mapping exists; mapping is
\code{ioremap}, a separate call, and a driver can map a range without claiming it
(\secref{c6:app:b2}). Absence means only that nobody has claimed the range. \code{qa-dev} was
present and answering at \code{0x0C000000} from the first boot, and Exercise~\ref{c1:ex:7} found
nothing there because no driver for it existed until this chapter's.

\xpart{d} The check. The code works because it still reaches the registers through \code{readl},
and C converts a \code{void *} to the accessor's \code{const volatile void __iomem *} argument
without complaint. What is gone is the annotation that made a mistake visible: the next change
that dereferences \code{regs} directly, or hands it to \code{memcpy}, compiles silently and passes
\code{sparse}, and on an architecture where registers need their own instructions or byte handling
that code is simply wrong. You have also lost the documentation: a reader can no longer tell from
the type that the pointer is not memory.

% ------------------------------------------------------------------------------------------
\solution[fillaccent2]{c6:ex:2}{The allocators}{Recall}

\xpart{a}
\begin{itemize}
  \item \textbf{48-byte state structure:} \code{kzalloc(sizeof(*priv), GFP_KERNEL)}. It is small,
    comes from a slab cache, and is zeroed, so no field starts as garbage.
  \item \textbf{2 MB software buffer:} \code{vmalloc}, or \code{kvmalloc}. Nothing needs it
    physically contiguous, and as one contiguous block it would be $2\,\mathrm{MB} / 4\,\mathrm{KB}
    = 512$ pages, order 9, which a system that has been up for a while often cannot supply.
  \item \textbf{16 KB DMA buffer:} \code{kmalloc(..., GFP_KERNEL)}, which is physically contiguous
    and at four pages, order 2, small enough to succeed on a fragmented system, then handed to the
    device through the DMA API with \code{dma_map_single}; or \code{dma_alloc_coherent}, which
    allocates and maps in one call. Never \code{vmalloc} (\secref{c6:app:a6}).
  \item \textbf{1 MB buffer hardware never sees:} \code{kvmalloc}. Contiguity is not needed, so it
    can try \code{kmalloc} for the 256 pages, order 8, and fall back to \code{vmalloc} when that
    fails. Free it with \code{kvfree}.
\end{itemize}

\xpart{b} 131072 bytes is 128 KB, $131072 / 4096 = 32$ pages, order 5. \code{kmalloc} serves
anything above 8 KB, its largest slab cache on 4 KB pages, straight from the page allocator, so the
request needs one free, physically contiguous block of 32 pages. On a freshly booted machine free
memory is in large blocks; after three weeks it is scattered in small ones, and there can be plenty
of it free with no block of order 5 anywhere. Order 5 is also above \code{PAGE_ALLOC_COSTLY_ORDER},
which is 3 (\file{include/linux/mmzone.h}): the allocator gives an allocation that size a limited
attempt at reclaim and compaction and then returns \code{NULL}, rather than retrying as it would for
a small one. The file is \file{/proc/buddyinfo}, which lists the free blocks of each order in each
zone; on the customer's machine the columns from order 5 up are zero or close to it, while
\file{/proc/meminfo} still shows memory free.

\xpart{c} Because you do not know which you got. \code{kvmalloc} returns \code{kmalloc} memory when
the \code{kmalloc} attempt succeeds and \code{vmalloc} memory when it does not, and nothing tells the
caller which. In testing the first attempt always succeeds, so the driver works; on the day
fragmentation defeats it, the buffer is \code{vmalloc} memory, which is not physically contiguous,
and \code{dma_map_single} refuses it with a warning and \code{DMA_MAPPING_ERROR}
(\code{dma_map_single_attrs} checks \code{is_vmalloc_addr}, \file{include/linux/dma-mapping.h}).
It fails exactly where fragmentation made the difference, which is in the field.

\xpart{d} \code{kmalloc} and \code{kzalloc} go with \code{kfree}, \code{vmalloc} with \code{vfree},
\code{kvmalloc} with \code{kvfree}, and \code{alloc_pages} with \code{__free_pages}, which takes the
\code{struct page *} and the order; \code{free_pages} takes an address and goes with
\code{__get_free_pages}. Mismatched, \code{vfree} of \code{kmalloc} memory finds no \code{vmalloc}
area, warns ``Trying to vfree() nonexistent vm area'' and frees nothing, which is a leak;
\code{kfree} of \code{vmalloc} memory takes an address outside the linear map for a slab object,
which is undefined and at best an immediate oops. The mismatch that survives testing is \code{kfree}
of \code{kvmalloc} memory: correct for as long as \code{kvmalloc} returns \code{kmalloc} memory,
which in testing is always, and a crash the first time it falls back. \code{kvfree} checks which it
was given, and is correct for both.

% ------------------------------------------------------------------------------------------
\solution[fillaccent3]{c6:ex:3}{GFP flags}{Hand calculation}

\xpart{a} \code{GFP_KERNEL}. \code{read()} runs in process context on behalf of the caller, where
sleeping is allowed, so nothing goes wrong, provided no spinlock is held at the allocation.

\xpart{b} \code{GFP_ATOMIC}. \code{GFP_KERNEL} may enter direct reclaim and sleep, and a handler may
not sleep. With \code{CONFIG_DEBUG_ATOMIC_SLEEP} the kernel prints ``BUG: sleeping function called
from invalid context'' on the first call, before it has needed to sleep; without it the code works
until the day the allocator has to reclaim, and then sleeps in interrupt context: ``BUG: scheduling
while atomic'' and quite possibly a hung CPU.

\xpart{c} \code{GFP_ATOMIC}; or better, allocate with \code{GFP_KERNEL} before taking the lock and
do only the linking-in under it. With \code{GFP_KERNEL} the holder may sleep with the lock held and
preemption disabled: ``scheduling while atomic'', every other CPU that wants the lock spinning while
it sleeps, and a deadlock if the task scheduled in its place on the same CPU wants the lock too.

\xpart{d} \code{GFP_KERNEL}. A module's init function runs in process context, in the
\code{finit_module} system call made by \code{insmod}, and may sleep. Nothing goes wrong.

\xpart{e} The only sense in which it is high priority is that it may dip into a reserve kept below
the normal watermark: \code{GFP_ATOMIC} is \code{__GFP_HIGH | __GFP_KSWAPD_RECLAIM}
(\file{include/linux/gfp_types.h}). What it cannot do is what makes an allocation succeed under
pressure: reclaim, write back, compact, and wait for any of them; it may only wake \code{kswapd}. So
\code{GFP_ATOMIC} is the one more likely to return \code{NULL}. When the free pages and the reserve
are gone it fails at once, where \code{GFP_KERNEL} would have reclaimed memory and succeeded. That
is why code using it must handle failure.

\xpart{f} No. It is a bug that has not fired yet. It works because every allocation so far found
memory without needing to reclaim, so it never slept. It appears when memory pressure forces the
allocator into direct reclaim while the spinlock is held: the task sleeps in atomic context, which
is ``scheduling while atomic'', a stall on every CPU waiting for the lock, or a deadlock.
\code{CONFIG_DEBUG_ATOMIC_SLEEP} would have reported it on the first run: every \code{kmalloc}
calls \code{might_alloc()}, which calls \code{might_sleep_if()} for any flag that allows blocking
(\file{include/linux/sched/mm.h}), so the check fires whenever the line runs, whether or not it
needed to sleep.

% ------------------------------------------------------------------------------------------
\solution[fillaccent3]{c6:ex:4}{Translating a device tree address}{Hand calculation}

\xpart{a} \code{0x40002000}, for \code{0x100} bytes, which is 256. The \code{reg} is read in the
bus's cells, one address and one size: child address \code{0x2000}, length \code{0x100}. The bus's
\code{ranges} maps child \code{0x0} to parent \code{0x0_40000000} for \code{0x1000000} bytes, and
\code{0x2000} lies inside that child window, so
\[
  \texttt{0x40000000} + (\texttt{0x2000} - \texttt{0x0}) = \texttt{0x40002000}.
\]
The parent is the root, whose address space is the CPU's physical address space, so the translation
ends there.

\xpart{b} Two: one address cell and one size cell. You know from the \emph{parent}: \code{reg} is
read in the \code{\#address-cells} and \code{\#size-cells} of the node's parent, and
\code{bus@40000000} declares one of each. The widget's own properties say nothing about it.

\xpart{c} Four. An entry in \code{ranges} is a child address, a parent address and a length. The
child address is in the bus's own \code{\#address-cells}, one cell: \code{0x0}. The parent address is
in the parent's, and the root declares two: \code{0x0 0x40000000}, the high cell and then the low,
so \code{0x0_40000000}. The length is in the bus's \code{\#size-cells}, one cell: \code{0x1000000}.
$1 + 2 + 1 = 4$.

\xpart{d} They map physical \code{0x2000}, which is not the widget but whatever this SoC decodes at
that address: on many parts boot ROM or on-chip SRAM, on some nothing at all. Most likely they read
something plausible, the contents of a ROM or zeros, with no error of any kind, as \code{base=0xC}
does on \code{qa-dev} in \secref{c6:app:b6}; if nothing decodes there, the first read raises a
synchronous external abort. Either way it is the identity check that tells them.

\xpart{e} Child address \code{0x0} on this bus is address \code{0x40000000} on the parent bus, for
\code{0x1000000} bytes, which is 16 MB: the bus's addresses \code{0x0} to \code{0xFFFFFF} are the
parent's \code{0x40000000} to \code{0x40FFFFFF}.

% ------------------------------------------------------------------------------------------
\solution[fillaccent4]{c6:ex:5}{Ordering}{Design}

\xpart{a} \code{writel} for \code{CTRL} and \code{readl} for \code{STATUS}, the ordered defaults.
With \code{_relaxed} for both, their order does not change: both go to one device through the
mapping \code{ioremap} made, and accesses from one CPU to one device arrive in program order
whether relaxed or not (\secref{c6:app:b5}; \file{Documentation/memory-barriers.txt}). What goes
wrong is something else, and it goes wrong with the ordered forms too. A conversion takes time, and
a status read issued straight after the command will usually find it not finished. The driver must
poll \code{STATUS} until it says done, with a timeout, which is what \code{readl_poll_timeout} does,
or wait for the device's interrupt.

\xpart{b} Nowhere among the register writes themselves. All seventeen go to one device, and relaxed
accesses to one device stay in program order, so the sixteen configuration writes land before
\code{GO} without any barrier. A barrier is needed in one place: before the \code{GO} write, and
only if the device will read something from RAM that the driver wrote first, a buffer or a
descriptor. So the sixteen are \code{writel_relaxed} and \code{GO} is a \code{writel}. That buys
sixteen barriers fewer on every start; on arm64 each \code{writel} puts a \code{dmb oshst} in front
of its store (\file{arch/arm64/include/asm/io.h} and \file{barrier.h}). It is a small saving, and
worth having only on a path where it has been measured to matter.

\xpart{c} Because what you observe points away from the cause. The failure is intermittent and
depends on timing, so it appears on one machine and not another, and under load rather than on the
bench. Adding a \code{printk} or attaching a debugger changes the timing and makes it vanish. It
shows up later and somewhere else, as data that is wrong, never at the line that is missing a
barrier. And the code, read line by line, is correct in program order, which is the order everyone
reads it in. A strongly ordered machine, such as an x86 desktop, never reorders a store to memory
past a later write to a register, so the code works there every time.

\xpart{d} Because real silicon is not lenient in any predictable way. On a real device a byte or
halfword write to a 32-bit register may be widened by the interconnect with the other bytes written
as zeros, or dropped, or answered with a bus error; a narrow read of a register with side effects,
such as \code{SAMPLE}, pops a whole sample and delivers a quarter of it. None of those is what the
driver meant, and which one you get depends on the chip. A model that widened quietly would let a
wrong driver pass every test here and fail on the first board; one that rejects the access makes
the bug fail here, where a failure costs a reboot. The documented access width is part of a
register's contract, and the model enforces it.

% ------------------------------------------------------------------------------------------
\solution[fillaccent4]{c6:ex:6}{Coherent or streaming}{Design}

\xpart{a} \code{dma_alloc_coherent}. Both sides touch small entries constantly: the driver writes
an entry, the device reads it and writes a status back, thousands of times a second. With a
streaming mapping every exchange would need a sync in each direction; coherent memory makes each
side's writes visible to the other without one, and the ring is small enough that slower uncached
CPU access costs nothing that matters. It still needs the \code{writel} barrier before the doorbell,
because coherent does not mean ordered (\secref{c6:app:a6}).

\xpart{b} \code{dma_map_single} with \code{DMA_FROM_DEVICE}, on a buffer from \code{kmalloc} (64 KB
is 16 pages, order 4) or \code{alloc_pages}. The device writes it once, and then the CPU reads all
64 KB in bulk to copy them to userspace. One mapping and one cache invalidate at the unmap, and the
CPU reads cached memory at full speed; in coherent memory, which is usually uncached, the 64 KB copy
would be slow.

\xpart{c} \code{dma_alloc_coherent}. Four bytes that the driver writes and the device reads over
and over: there is no bulk access to make faster, and a streaming mapping would need a
\code{dma_sync_single_for_device} after every write. It is exactly the kind of status word
\secref{c6:app:a6} names.

\xpart{d} It broke the ownership rule: between the map and the unmap the buffer belongs to the
device, and the CPU read it without borrowing it back. On a machine whose DMA is not coherent with
the CPU's caches, the read pulls a line of the buffer into the cache holding whatever was in memory
before the device wrote it, so the peek sees stale data; depending on what the architecture does at
the unmap, that stale line can also outlive the transfer, and the driver then reads old data where
the device's should be. Where the buffer is bounced, because the device cannot reach the memory it
is in, the device writes a bounce buffer and the peek reads the original, which the kernel fills
only at the unmap. Testing did not catch it because the machines it ran on were coherent, an x86
desktop, an arm64 system with coherent DMA, or an emulator with no caches to go stale, and nothing
was bounced; there the read is harmless. It breaks on an SoC whose DMA master is not cache-coherent,
which is common in embedded parts, and on any system that bounces the buffer. The fix is
\code{dma_sync_single_for_cpu} before the peek and \code{dma_sync_single_for_device} after it.

\xpart{e} \code{kvmalloc} is exactly wrong here. On the fragmented system they are worried about,
the \code{kmalloc} attempt fails and \code{kvmalloc} falls back to \code{vmalloc}, whose memory is
not physically contiguous and cannot be handed to a device. The kernel refuses it:
\code{dma_map_single_attrs} checks \code{is_vmalloc_addr} and returns \code{DMA_MAPPING_ERROR}
with the warning ``rejecting DMA map of vmalloc memory'' (\file{include/linux/dma-mapping.h}). So
the driver works in testing and fails in the field, on precisely the system the change was meant to
help. Use \code{kmalloc}, allocated once when the device is opened or probed and kept rather than
per transfer, with its failure handled; or \code{alloc_pages} or \code{dma_alloc_coherent}; or, if
the device can do scatter-gather, a list of single pages mapped with \code{dma_map_sg}, which needs
no contiguity at all.

% ------------------------------------------------------------------------------------------
\solution[fillaccent4]{c6:ex:9}{What a wrong address does}{Design}

\xpart{a} Zeros: the course measured \code{id=0x00000000 version=0x0000} (\secref{c6:app:b6}), and
nothing faults, because \code{0xC} lies in the virt machine's flash window, which starts at
physical 0 and is \code{0x08000000} bytes long (\file{hw/arm/virt.c} in QEMU), and reads as zeros
here. The identity check compares \code{0x00000000} with \code{0x51414456}, fails, unwinds,
\code{iounmap} and then \code{release_mem_region}, and returns \code{-ENODEV}; \code{insmod} reports
``No such device'' and nothing is left behind. A driver without the check would have carried on as
if it had the device: written \code{0xDEADBEEF} to \code{SCRATCH}, which here is physical
$\texttt{0xC} + \texttt{0x30} = \texttt{0x3C}$, into whatever lives there, and then gone on to
configure and read a device that is not there, failing later and elsewhere with nothing in the log
pointing at the address.

\xpart{b} What the course recorded (\secref{c6:app:b6}): the read of \code{QA_DEV_ID} at
\code{0xC001000} raises ``Internal error: synchronous external abort: 0000000096000010'', with the
program counter in \code{qa_mmio_init+0x98}, called from \code{do_one_initcall}. The syndrome
decodes as a data abort taken in the kernel (the top six bits, \code{0x25}) on a read, with fault
status \code{0x10}, a synchronous external abort. Nothing decodes that address: the platform bus
window is \code{0x0C000000} to \code{0x0DFFFFFF}, and \code{qa-dev} occupies only its first page,
\code{0x0C000000} to \code{0x0C000FFF}. \code{insmod} is killed. The module is left half-loaded,
its init never having returned, so it can be neither used nor removed, and it still holds the
region it claimed. The only way out is a reboot.

\xpart{c} ``It did not crash'' is not evidence that a mapping is right. A wrong address can read
plausible values with no error, or abort the machine, depending only on what the machine decodes
there, and the quiet outcome is the more dangerous of the two. The evidence has to be positive: an
identity register that reads the value the datasheet promises, and a scratch register that reads
back what was written.

\xpart{d} On a board on your bench, the wrong-but-mapped address may belong to another peripheral,
so the scratch write can reprogram something real, and the abort costs a power cycle and a walk
across the room, or a JTAG probe if it happened early in boot; on a product in the field, either
one is a device that fails or hangs at a customer's, and bringing it back costs a truck or a recall.
\Secref{c1:app:a4} named what the course gives up by booting in QEMU, the experience of real boards
and of recovering one, and this is the other side of that trade: here both experiments are free,
repeatable and two seconds long, which is why the course asks you to run them.

% ------------------------------------------------------------------------------------------
\solution[fillaccent]{c6:ex:10}{An address worked out by hand}{Cross-check}

\xpart{a} \code{reg} is eight bytes, \code{00 00 00 00 00 00 10 00}, which is two cells,
\code{0x00000000} and \code{0x00001000}. \code{ranges} is sixteen bytes,
\code{00 00 00 00 00 00 00 00 0c 00 00 00 02 00 00 00}, which is four cells: \code{0x00000000},
\code{0x00000000}, \code{0x0C000000} and \code{0x02000000}. \code{reg} has two cells because the
platform bus declares one address cell and one size cell; \code{ranges} has four because its
parent address is in the root's two cells. QEMU writes exactly these: \code{ranges} as 0, the high
and low halves of the bus base, and the bus size (\file{hw/core/sysbus-fdt.c}), and \code{reg} as
the device's offset on the bus and \code{QA_DEV_MMIO_SIZE} (\file{tools/qa-dev/integrate.py}).

\xpart{b} The two numbers to add are \code{reg}'s address cell and the parent address in
\code{ranges}, not either size:
\[
  \texttt{0x0C000000} + (\texttt{0x00000000} - \texttt{0x00000000}) = \texttt{0x0C000000},
\]
the child's address less the child base of the range, added to the parent base. The high parent
cell is 0, so the full 64-bit address is \code{0x00000000_0c000000}. The window is \code{0x1000}
bytes, 4096, one page, which is \code{QA_DEV_MMIO_SIZE}; it ends at \code{0x0C000FFF}, well inside
the bus window, which runs to $\texttt{0x0C000000} + \texttt{0x02000000} - 1 =
\texttt{0x0DFFFFFF}$.

\xpart{c} Offset 0 is \code{QA_DEV_ID}, and the header gives its value: \code{QA_DEV_ID_MAGIC},
\code{0x51414456}, which is ``QADV'' in ASCII, big end first (\code{0x51} Q, \code{0x41} A,
\code{0x44} D, \code{0x56} V). Offset 4 is \code{QA_DEV_VERSION}, for which the header gives only
the layout, major in bits 15:8 and minor in bits 7:0; but the compatible string it defines,
\code{qacademy,qa-dev-1.0}, says version 1.0, so predict \code{0x00000100}.

\xpart{d} A correct module at that address prints \code{id=0x51414456 version=0x0100}; those are
the two strings \code{make test L=L06} looks for in the kernel log.

\xpart{e}
\begin{itemize}
  \item It should match. \file{run.sh} computes the same address the same way, prints ``so the
    registers are at 0xC000000, 4096 bytes'', and loads the module with \code{base=201326592},
    which is \code{0x0C000000} in decimal. If yours differed, you took \code{reg} at face value
    (\code{0x0}), added the wrong cell of \code{ranges}, the high half of the parent address or the
    size, or read the bytes with \code{-tx4}. \file{run.sh} reads raw bytes with \code{od -tx1},
    reassembles them four at a time, most significant first, and adds \code{reg}'s first cell to
    the third cell of \code{ranges}.
  \item \code{-tx1} is right. \code{-tx4} decodes each group of four bytes as a native-endian
    integer, and the target is little-endian, so every big-endian cell comes back byte-swapped:
    \code{reg} reads \code{00000000 00100000} and \code{ranges} reads \code{00000000 00000000
    0000000c 00000002}. The same arithmetic then gives an address of \code{0xC} and a window of
    \code{0x100000}, a round megabyte. It is plausible because it has the right shape: small hex
    numbers in the right places, and an address that is a real address, which reads zeros without
    complaint (\secref{c6:app:b6}).
  \item The constant is \code{0x0C000000}. It is the parent base in \code{ranges}, declared by the
    parent node, \code{platform-bus@c000000}; it comes from QEMU's virt machine, which places its
    platform bus at \code{0x0C000000} with \code{0x02000000} bytes (\file{hw/arm/virt.c}) and
    writes that into \code{ranges} when it generates the device tree.
  \item It is in the device tree QEMU generated and handed to the kernel, readable under
    \file{/sys/firmware/devicetree/base}, and today on \code{insmod}'s command line, where
    \file{run.sh} put it. For the driver to find out, it has to become a platform driver matching
    the compatible \code{qacademy,qa-dev-1.0}: the kernel then creates the device from the node,
    translates its \code{reg} through every \code{ranges} above it into a \code{struct resource},
    and hands it to \code{probe}, where \code{platform_get_resource}, or
    \code{devm_platform_ioremap_resource}, gives the driver its registers. That is
    Chapter~\ref{c10:ch}.
\end{itemize}

\xpart{f} $\texttt{0x0C000000} + (\texttt{0x1000} - \texttt{0x0}) = \texttt{0x0C001000}$, and
\code{QA_DEV_ID} still reads \code{0x51414456}. The node is not a separate claim about the device;
it is QEMU's description of where it put it, generated from the address it assigned
(\file{tools/qa-dev/integrate.py}). A node that says \code{0x1000} means the device sits one page
further up, and its identity register is at offset 0 of its window wherever that is: the identity
belongs to the device, not to the address. Only if the node were edited by hand while the device
stayed at \code{0x0C000000} would \code{0x0C001000} be the unbacked page of \secref{c6:app:b6}, and
the read abort.
